%% 
%% Copyright 2007-2025 Elsevier Ltd
%% 
%% This file is part of the 'Elsarticle Bundle'.
%% ---------------------------------------------
%% 
%% It may be distributed under the conditions of the LaTeX Project Public
%% License, either version 1.3 of this license or (at your option) any
%% later version.  The latest version of this license is in
%%    http://www.latex-project.org/lppl.txt
%% and version 1.3 or later is part of all distributions of LaTeX
%% version 1999/12/01 or later.
%% 
%% The list of all files belonging to the 'Elsarticle Bundle' is
%% given in the file `manifest.txt'.
%% 
%% Template article for Elsevier's document class `elsarticle'
%% with numbered style bibliographic references
%% SP 2008/03/01
%% $Id: elsarticle-template-num.tex 272 2025-01-09 17:36:26Z rishi $
%%
\documentclass[preprint,12pt,authoryear]{elsarticle}

%% Use the option review to obtain double line spacing
%% \documentclass[authoryear,preprint,review,12pt]{elsarticle}

%% Use the options 1p,twocolumn; 3p; 3p,twocolumn; 5p; or 5p,twocolumn
%% for a journal layout:
%% \documentclass[final,1p,times]{elsarticle}
%% \documentclass[final,1p,times,twocolumn]{elsarticle}
%% \documentclass[final,3p,times]{elsarticle}
%% \documentclass[final,3p,times,twocolumn]{elsarticle}
%% \documentclass[final,5p,times]{elsarticle}
%% \documentclass[final,5p,times,twocolumn]{elsarticle}

%% For including figures, graphicx.sty has been loaded in
%% elsarticle.cls. If you prefer to use the old commands
%% please give \usepackage{epsfig}

%% The amssymb package provides various useful mathematical symbols
\usepackage{amssymb}
%% The amsmath package provides various useful equation environments.
\usepackage{amsmath}
\usepackage{subfiles}
\usepackage{algorithm}
\usepackage{algorithmic}
\usepackage{xurl}
\usepackage{amssymb}
\usepackage{pifont}
\usepackage{booktabs}
\usepackage{longtable}

% Define custom checkmarks and crosses
\newcommand{\cmark}{\ding{51}}  % ✓ checkmark
\newcommand{\xmark}{\ding{55}}  % ✗ cross mark
\newcommand{\hmark}{$\circ$\hspace{-0.66em}$\bullet$}  % ◐ half-filled circle

\usepackage{hyperref}
\hypersetup{breaklinks=true, colorlinks=true, urlcolor=blue}
%% The amsthm package provides extended theorem environments
%% \usepackage{amsthm}

%% The lineno packages adds line numbers. Start line numbering with
%% \begin{linenumbers}, end it with \end{linenumbers}. Or switch it on
%% for the whole article with \linenumbers.
%% \usepackage{lineno}

\journal{Nuclear Physics B}

\begin{document}

\begin{frontmatter}

%% Title, authors and addresses

%% use the tnoteref command within \title for footnotes;
%% use the tnotetext command for theassociated footnote;
%% use the fnref command within \author or \affiliation for footnotes;
%% use the fntext command for theassociated footnote;
%% use the corref command within \author for corresponding author footnotes;
%% use the cortext command for theassociated footnote;
%% use the ead command for the email address,
%% and the form \ead[url] for the home page:
%% \title{Title\tnoteref{label1}}
%% \tnotetext[label1]{}
%% \author{Name\corref{cor1}\fnref{label2}}
%% \ead{email address}
%% \ead[url]{home page}
%% \fntext[label2]{}
%% \cortext[cor1]{}
%% \affiliation{organization={},
%%             addressline={},
%%             city={},
%%             postcode={},
%%             state={},
%%             country={}}
%% \fntext[label3]{}

\title{An Interpretable Temporal Graph Learning Method for Real-Time Financial Fraud Detection and Decision Making}


% \author{Quang-Vinh Dang}
% \ead{vinh.dq4@buv.edu.vn}
% \affiliation{organization={School of Innovative and Computing Technologies, British University Vietnam},
%             city={Hung Yen},
%             postcode={160000},
%             country={Vietnam}}

% \author{Ngoc-Son-An Nguyen}
% \affiliation{organization={Faculty of Information Technology, Industrial University of Ho Chi Minh City},
%             city={Ho Chi Minh City},
%             postcode={700000},
%             country={Vietnam}}

% \author{Thi-Bich-Diem Vo}
% \affiliation{organization={GiaoHangNhanh},
%             city={Ho Chi Minh City},
%             postcode={700000},
%             country={Vietnam}}

            
% \author{Minh Ngoc Dinh}
% \affiliation{organization={Millennia University},
%             city={Ho Chi Minh City},
%             postcode={700000},
%             country={Vietnam}}


%% Abstract
\begin{abstract}
Financial fraud detection systems face a critical challenge: while Graph Neural Networks (GNNs) have demonstrated strong predictive performance, they often rely on post hoc explanation techniques that introduce substantial computational overhead and limit their suitability for real-time, regulation-compliant decision support. This limitation is particularly important in environments subject to transparency requirements, such as the General Data Protection Regulation (GDPR) Article 22 and the Equal Credit Opportunity Act (ECOA). To address this challenge, this study proposes RTXGNN, an intrinsically interpretable Graph Neural Network framework that jointly optimizes fraud detection performance and real-time explainability without requiring post hoc explanation generation. RTXGNN incorporates two key components: (1) Hierarchical Recency-Aware Positional Encoding (HRAPE), which captures multi-scale temporal characteristics including recency, periodicity, and transaction velocity; and (2) a Self-Explainable Aggregation Layer (SEAL), which generates multi-granularity explanations during the inference process. The proposed framework is evaluated on the Elliptic Bitcoin transaction dataset, which contains 203,769 transactions across 49 temporal steps. Experimental results demonstrate strong predictive performance, achieving an F1-score of 0.5129 and an area under the receiver operating characteristic curve (AUC) of 0.8661, representing a 19.3\% improvement in F1-score over the strongest baseline method. Ablation analyses show that HRAPE contributes a 2.53 percentage-point improvement in F1-score, while sparsity regularization yields an additional 6.37 percentage-point gain. Furthermore, RTXGNN preserves 76\% of its full-supervision performance when trained with only 10\% labeled data and achieves positive explanation fidelity (0.0764). The results suggest that explainability-oriented design can serve not only as a transparency mechanism but also as a form of domain-informed regularization, enhancing both predictive performance and interpretability in financial fraud detection.
\end{abstract}

% %%Graphical abstract
% \begin{graphicalabstract}
% %\includegraphics{grabs}
% \end{graphicalabstract}

%%Research highlights
% \begin{highlights}
% \item Research highlight 1
% \item Research highlight 2
% \end{highlights}

\begin{keyword}
Financial fraud analytics; Explainable analytics; Graph learning; Decision support; Temporal networks; Regulatory compliance
\end{keyword}

%% PACS codes here, in the form: \PACS code \sep code

%% MSC codes here, in the form: \MSC code \sep code
%% or \MSC[2008] code \sep code (2000 is the default)


\end{frontmatter}

%% Add \usepackage{lineno} before \begin{document} and uncomment 
%% following line to enable line numbers
%% \linenumbers

%% main text
%%
\subfile{1introduction}
\subfile{2related}
\subfile{3model}
\subfile{4method}
\subfile{5eval}
\subfile{6conclusions}

\section{Declaration of generative AI and AI-assisted technologies in the writing process}

During the preparation of this work the author(s) used ChatGPT in order to polish the writing. After using this tool/service, the author(s) reviewed and edited the content as needed and take(s) full responsibility for the content of the published article.

%% If you have bib database file and want bibtex to generate the
%% bibitems, please use
%%
%%  \bibliographystyle{elsarticle-num} 
%%  \bibliography{<your bibdatabase>}

%% else use the following coding to input the bibitems directly in the
%% TeX file.

%% Refer following link for more details about bibliography and citations.
%% https://en.wikibooks.org/wiki/LaTeX/Bibliography_Management
% \bibliographystyle{elsarticle-num} 
\bibliographystyle{elsarticle-harv}
\bibliography{reference}

\end{document}

\endinput
%%
%% End of file `elsarticle-template-num.tex'.
